\subsection{Comparing equilibrium strategies and offsetting effects}

The analysis compares all ordered pairs changing one legal or preference dimension
within the core matrix. At each of five costs, the American, Trial Fee-Shifting,
and Complete Fee-Shifting rules give six directed fee changes within each risk
preference. Both directions of preference change are also computed within every
fee rule. There are 18 comparisons per cost and 90 in total. Both parties'
preferences change in a preference comparison. Reverse comparisons are calculated
independently, not obtained by negating forward contributions. The tables report strategies from the original computed equilibria,
not the more broadly mixed profiles used in the sensitivity analysis.

\paragraph{Information sets and measurement.}
Strategies are compared at information sets reached with positive probability
in both endpoint equilibria. Each row conditions on the acting party's own
signal and observed history. The signal is on the common plaintiff-case-strength
scale, so a higher signal is more favorable to the plaintiff and less favorable
to the defendant. Offer information sets additionally distinguish the party's
private commitment to continue or exit if bargaining fails. A commitment to
abandon or default is not an unconditional probability of actual abandonment or
default, because settlement can occur before the commitment takes effect.

Filing, answering and exit commitments are reported as probabilities, expressed
as percentages. Their changes and decomposed contributions are expressed in
percentage points (pp): a change from 87.8\% to 100\% is an increase of
12.2 percentage points, not a 12.2\% proportional increase. Pure demands and
offers, and their decomposed contributions, are expressed as fractions of the
normalized award. For example, a demand changing from 0.75 to 0.05 changes by
$-0.70$. Offer amounts are used only when the endpoint strategies and the
counterfactual strategies entering the accounting are pure. Mixed offers are
not replaced by mean offers in these tables.

\paragraph{Best-response comparisons.}
Let $\theta_0$ and $\theta_1$ denote the original and target fee-rule or preference
specifications, and let $\sigma^0$ and $\sigma^1$ denote their selected equilibrium
strategy profiles. For each player, the opponent's strategy is partitioned into
three components: entry (filing or answering), offers, and exit commitments.
Under $\theta_1$, all eight combinations of original and target opponent
components are constructed. Each comparison holds the opponent's complete
hybrid strategy fixed and computes an unrestricted best response for the focal
player, jointly reoptimizing that player's current and subsequent decisions.
The calculation does not impose monotonicity or simulate an observed
adjustment process.

Write $z(S)$ for the reported strategy coordinate in this best response when
the set $S$ of opponent components has been replaced by its target-equilibrium
values. The coordinate is the probability of the named binary action, a pure demand or
offer amount, or an action probability for a mixed offer. Write $x$ and $y$ for the corresponding original
and target equilibrium coordinates. The direct contribution is
\[
  D=z(\varnothing)-x.
\]
For each of the six orders of replacing the three opponent components, the
marginal change from replacing component $k$ is recorded. Its contribution is
the average of those increments:
\[
  C_k=\frac{1}{6}\sum_{\pi}
  \left[z\bigl(S_{\pi,k}\cup\{k\}\bigr)-z(S_{\pi,k})\right],
\]
where $S_{\pi,k}$ contains the components preceding $k$ in order $\pi$. Thus
\[
  y-x=D+C_{\mathrm{entry}}+C_{\mathrm{offers}}+C_{\mathrm{exit}}+R,
  \qquad R=y-z(\{\mathrm{entry},\mathrm{offers},\mathrm{exit}\}).
\]
This is a direct-first accounting convention. Interactions between the change
in rules or preferences and the opponent's adjustment enter the opponent
contributions. These quantities are not uniquely identified causal shares.
Individual contributions can exceed the net change or have the opposite sign
because other contributions offset them.

\paragraph{Selection of the displayed changes.}
The changed-action sections report a restricted subset of actual strategy changes. A changed
information set qualifies for the focus criterion if, against the target
opponent and with subsequent own decisions optimized, the original local
strategy assigns probability mass greater than $10^{-6}$ to actions whose
conditional utility is more than $10^{-6}$ below the best action. The utility
threshold is in the target specification's reported utility units. This is a
conditional action comparison, not the loss from reverting the entire strategy
or a welfare comparison across preference specifications. In particular,
disjoint strategy supports do not suffice: an action absent from the target
support can nevertheless remain optimal.

The criterion is evaluated for the original equilibrium profiles and for two
tested sets of equilibrium-preserving mixed representatives. The
changed-action sections retain its intersection across these three comparisons.
The verification inventory records selected-coordinate counts for the complete
90-comparison collection. Consecutive signals are grouped
when their original and target distributions and numerical contributions agree
within $10^{-6}$, even if their sensitivity flags differ. ``At some signals''
means that at least one, but not every, signal included in that grouped row has
a tie- or off-path-sensitive allocation.

The Remaining column always preserves the endpoint-selection residual, including
nonzero residuals. It is not allocated to another component. Undefined conditional
comparisons are displayed as dashes. Selection for this display does not imply that other
strategy changes are unimportant or explained. Changes involving newly reached
or no-longer-reached histories, and redistribution among actions that remain
optimal, are outside the restricted table.

\paragraph{Unchanged actions with offsetting effects.}
Restricting attention to changed actions can conceal counteracting incentives.
An intervention can make the original action unattractive against the original
opponent even though the opponent's equilibrium adjustment restores that action.
The additional sections therefore retain selected information sets where the
original and target local distributions agree within $10^{-6}$, but the direct
best response differs. Against the original opponent under the target rules or
preferences, the original local policy must assign probability mass greater
than $10^{-6}$ to actions more than $10^{-6}$ below the best conditional action,
with subsequent own decisions optimized. This excludes a direct switch caused
solely by selecting differently among tied or nearly tied actions. The best
response to the fully updated opponent must also reproduce the target local
distribution within $10^{-6}$. All eight conditional comparisons must be defined,
and the reported direct contribution must be nonzero and the remainder no
larger than $10^{-6}$ in the reported coordinate.

These requirements are applied separately to the original profiles and both
tested mixed representations, retaining their intersection. The publication
JSON distinguishes all selected coordinates, including unchanged-action rows.
An empty selection means no rows satisfy the implemented criteria; it does not
establish equality of the full strategies. Displayed endpoints and contributions
still come from the original profiles. Inclusion surviving these mixing checks
does not imply that the precise allocation is insensitive to ties or off-path
completions; the sensitivity column continues to report those checks.

For these rows $x=y$ and $R$ is numerically zero, so the direct contribution is
offset by the sum of the opponent contributions. For example, under the
American rule, the plaintiff files at signal 0.35 in both the risk-neutral and
moderately risk-averse equilibria. Changing preferences while retaining the
original opponent produces a best response of not filing: a direct contribution
of $-100$ percentage points. The opponent-offer contribution is $+100$
percentage points, with zero entry and exit contributions. This accounting
shows why an unchanged equilibrium action need not imply an absence of
counteracting incentives. It does not establish a uniquely necessary mechanism
or claim that the counterfactual actions occurred during equilibrium adjustment.

\paragraph{Equilibrium-preserving mixing checks.}
The auxiliary mixing search increases the mean normalized quadratic mixing
score $(1-\sum_a p_a^2)/(1-1/A)$ over information sets reached in the original
equilibrium, where $A$ is the full number of available actions. A player's
same-decision information sets are varied jointly. Candidate actions comprise
current support and actions tied in conditional utility under current
continuation play. Each proposed complete profile is checked against both
players' unrestricted best responses; a utility tie for the acting player alone
does not authorize adding an action to equilibrium support. Source-unvisited
policies remain fixed.

The main search admits maximum unilateral root-payoff gain of $10^{-9}$ reported
utility units and evaluates forward and reverse decision-block orders. The
higher-scoring verified result is selected. A second calculation uses forward
order and a tenfold tighter gain limit; its conditional tie threshold is also
tightened from $10^{-10}$ to $10^{-11}$. The search is local and does not certify
a globally maximally mixed or unique equilibrium. Every search is limited to
six sweeps; individual stopping records remain part of the supporting data.
The intersection therefore establishes stability only across the tested
representations. Even within that intersection, the decomposition can change
when a different mixed representative is used.

\paragraph{Ties, reach and verification.}
Primary best responses retain and renormalize the original probability on
actions that remain optimal within the $10^{-10}$ tie tolerance; when no such
probability remains, the first optimal action is selected. Alternative low- and
high-action selections test tie sensitivity. Separate checks replace opponent
policies unvisited in their donor equilibrium with extreme low- or high-action
completions. These are sensitivity checks, not bounds over all admissible
completions. The sensitivity column reports whether either check changes a
contribution in the original-profile accounting. It does not certify stability
under all equilibrium-preserving mixtures.

An intermediate best response may not actually reach an information set that is
reached in both endpoint equilibria. Conditional continuation comparisons are
still defined when chance-and-opponent reach is positive; the focal player's
own prior action probabilities are omitted from that reach calculation. If
chance-and-opponent reach is zero, no conditional allocation is imputed. The
underlying diagnostics retain these reach distinctions.
The table marks a row with an asterisk if at least one of its primary hybrid
best responses does not actually reach the information set (for at least one
signal in a grouped row). Its conditional continuation comparison remains
defined by positive chance-and-opponent reach; it is not an action predicted
to occur on that intermediate best-response path. This marker is distinct from
the tie/off-path sensitivity column. Complete best responses
are independently replayed to check root utilities and conditional action
values, and the saved equilibrium profiles and reports are checked against the
initialized game. Reported table entries are rounded; accounting identities are
verified before rounding.
